%
% Auteur : Olivier Cots
% Date   : 17 avril 2019
%
% macros latex
%


\newcommand*\myurl[1]{\href{#1}{#1}} % because of a bug using url after homotopy

% for homotopy
\def\sb{\bar{s}}
\def\cb{\bar{c}}
\def\lb{\bar{\lambda}}
\def\sphere{\mathbb{S}}
\def\laz{\lambda_0}
% end for homotopy


\usepackage[strict]{changepage}

% \mnote[vertical shift]{l/r}{text} 
\newcommand*\mnote[3][0pt]{%

    \checkoddpage

    \if l#2
        \ifoddpage\reversemarginpar\else\normalmarginpar\fi
        \def\pointer{\blacktriangleright}%
        \def\stackalignment{r}
        \def\widthbox{2.1cm}
        \def\offsetbox{1.6cm}
    \fi%
    \if r#2
        \ifoddpage\normalmarginpar\else\reversemarginpar\fi
        \def\pointer{\blacktriangleleft}%
        \def\stackalignment{l}
        \def\widthbox{2.1cm}
        \def\offsetbox{0.1cm}
    \fi%
    \marginpar{%
        \hspace*{\offsetbox}%
        \topinset{%
        \scalebox{1.5}{\small\textcolor{blue}{$\pointer$}}}{%
        \belowbaseline[-1.5\baselineskip-#1]{%
            \stackengine%
            {-5pt}%
            {\fcolorbox{blue}{white}{\parbox{\widthbox}%
                {\footnotesize\vspace{3pt}\raggedright\textcolor{black}{#3}\vspace{3pt}}}}%
            {~\colorbox{white}{\sffamily \textcolor{blue}{Note}}}%
            {O}%
            {l}%
            {F}%
            {F}%
            {S}%
            }%
        }{%
        #1+3ex}{%
        -2.2ex}%
    }%
}

\newcommand*\mynote[1]{\mnote[0pt]{l}{#1}}

%---------------------------------------------------------------------------------------------------------------
% Repertoire contenant les videos
\def\repVideos{Videos/}
%\def\repVideos{./}

%---------------------------------------------------------------------------------------------------------------
% Mise en valeur
%\newcommand{\myemph}[1]{\emph{\textbf{#1}}}                 % remplace \emph
%\newcommand{\empha}[1]{\textbf{\textcolor{blue}{#1}}}
%\newcommand{\emphb}[1]{\textcolor{red}{#1}}
\newcommand{\cblue}[1]{{\color{blue} #1}}
\newcommand{\cred}[1]{{\color{red} #1}}
\newcommand{\cgreen}[1]{{\color{green} #1}}
\newcommand{\crose}[1]{{\color{rose} #1}}
\newcommand{\clink}[1]{{\color{linkcol} #1}}
\def\emphb{\cblue}
\def\emphr{\cred}
\def\emphg{\cgreen}
\def\emphp{\crose}
\def\emphl{\clink}

%
\newcommand{\myemph}[1]{\textbf{\textcolor{blue}{#1}}}

% Cross mark and check mark
\newcommand{\cmark}{\ding{51}}%
\newcommand{\xmark}{\ding{55}}%

% item
%\newcommand{\newpoint}{{\scriptsize\emphb{\raisebox{1.3pt}{$\blacktriangleright$}~}}}
\newcommand{\newpoint}{\emphb{\raisebox{0pt}{\textbullet}~}}

\newcommand{\retourligne}{~\\ \vspace{-\baselineskip}}

\ifSLIDES

    %---------------------------------------------------------------------------------------------------------------
    %---------------------------------------------------------------------------------------------------------------
    % Beamer
    
    %---------------------------------------------------------------------------------------------------------------
    % Bullets
    \newcommand{\myitem}[1]{\vfill #1}
    \newcommand{\myitembulletindent}{\vfill \hspace{1em} {\small \textbullet} \hspace{0.1em} }
    \newcommand{\myitembullet}{{\small  \textbullet} }
    \newcommand{\myitemtiret} {\hspace{1em} -- }
    \newcommand{\myitemtiretnoi} {\vfill -- }

\else

    \newcommand{\myitembullet}{\noindent{\small  \textbullet} }

\fi


%---------------------------------------------------------------------------------------------------------------
%  Commande pour creer un nouveau slide
%  Couleur cfond doit etre definie

% myframe beamer
%\makeatletter

\def\colorbeamerheading{cfondtitre}
\newcommand{\setcolorbeamerheading}[1]{
    \def\colorbeamerheading{#1}
}

\newenvironment{myframe}[2][]
{
    \ifx\relax#1\relax%
    \else
        \def\colorbeamerheading{#1} % ce ne change que localement
    \fi
    \begin{frame}%
        \begin{minipage}{0.9\textwidth} 
            \large \color{\colorbeamerheading}\textsc{#2}
        \end{minipage}
        \hfill%
        \begin{minipage}{0.04\textwidth}
            \begin{flushright}
                \compteurPage[\colorbeamerheading]
            \end{flushright}
        \end{minipage}
        \vspace{-2pt}
        {\color{\colorbeamerheading}\noi\rule{\textwidth}{0.1em}}%
}
{%
    \end{frame}
}


%\makeatother

%\def\sizefontslide{\normalsize}
%\newcommand{\myframeold}[2][cfond]{
%    \addtocounter{pages}{1}
%    \vfill\null\end{slide}\\
%    \begin{slide} 
%        \hskip -.02\paperwidth  
%        %\MyColorBox[#1]{
%            \begin{minipage}{1.0\paperwidth}
%                \vspace{1em} 
%                \begin{minipage}{0.9\paperwidth} 
%                    \hspace{0.24in} \large \color{#1}\textsc{#2}
%                \end{minipage}
%                \hfill
%                \begin{minipage}{0.09\paperwidth}
%                    \compteurPage[#1]
%                \end{minipage}
%
%                {\color{#1}\noi\rule{\linewidth}{0.2em}}
%            \end{minipage}
%        %}
%        \vspace{0.1in}\sizefontslide
%}
%%\newcommand{\newslide}{\frame}

\newcommand{\frameNoTitle}{
    \addtocounter{pages}{1}
    \vfill\null\end{slide}\\
    \begin{slide} 
        \hskip -.02\paperwidth  

        \vspace{0.1in}\sizefontslide
}

%---------------------------------------------------------------------------------------------------------------
% Softwares, etc
\newcommand{\vrb}[1]{\texttt{#1}}
\newcommand{\cmd}[1]{\texttt{#1}}
\newcommand{\lib}[1]{\texttt{#1}}
\def\hampath{\lib{HamPath}}
\def\bocop{\lib{Bocop}}
\newcommand{\matlab}{\lib{Matlab}}
\newcommand{\octave}{\lib{Octave}}
\newcommand{\fortran}{\lib{Fortran}}
\newcommand{\python}{\lib{Python}}
\newcommand{\interface}{\lib{Interface}}
\newcommand\tapenade{\lib{TAPENADE}}

%---------------------------------------------------------------------------------------------------------------

%---------------------------------------------------------------------------------------------------------------
% Homepage
%\newcommand{\homepage}{http://cots.perso.enseeiht.fr}
%\newcommand{\resources}{\homepage/resources}

%---------------------------------------------------------------------------------------------------------------
% Misc
\newcommand{\noi}{\noindent}
\newcommand{\HRule}{\noi\rule{\linewidth}{0.3mm}}

%---------------------------------------------------------------------------------------------------------------
% Tag a gauche ou a droite : utile pour la numerotation des equations
\makeatletter
\newcommand{\reqnomode}{\tagsleft@false\let\veqno\@@eqno}
\newcommand{\leqnomode}{\tagsleft@true\let\veqno\@@leqno}
\makeatother

%---------------------------------------------------------------------------------------------------------------
% Litteraire
\newcommand\ie{\textit{i.e.}}
\newcommand\cf{{cf.}}

%---------------------------------------------------------------------------------------------------------------
% Notes personnelles : ne sont visibles que si \VIEWALLtrue est present dans le main
\newcommand{\notePerso}[1]{\ifVIEWALL{ {\color{green} NotePerso~:~#1} }\else\fi}
\newcommand{\noteInMarge}[1]{\ifVIEWALL{ \marginpar{{\color{red} #1}} }\else\fi}        % Pour avoir des notes dans la marge
\newcommand{\todo}[1]{\ifVIEWALL{\noindent {\color{red} \it TODO!!! #1}~\\}\else\fi}    % Pour ecrire des todoliste

%---------------------------------------------------------------------------------------------------------------
% Remplace la commande ``bar'' par ``xoverline''
%
\makeatletter
\newsavebox\myboxA
\newsavebox\myboxB
\newlength\mylenA

\newcommand*\xoverline[2][0.75]{%
    \sbox{\myboxA}{$\m@th#2$}%
    \setbox\myboxB\null% Phantom box
    \ht\myboxB=\ht\myboxA%
    \dp\myboxB=\dp\myboxA%
    \wd\myboxB=#1\wd\myboxA% Scale phantom
    \sbox\myboxB{$\m@th\overline{\copy\myboxB}$}%  Overlined phantom
    \setlength\mylenA{\the\wd\myboxA}%   calc width diff
    \addtolength\mylenA{-\the\wd\myboxB}%
    \ifdim\wd\myboxB<\wd\myboxA%
       \rlap{\hskip 0.5\mylenA\usebox\myboxB}{\usebox\myboxA}%
    \else
        \hskip -0.5\mylenA\rlap{\usebox\myboxA}{\hskip 0.5\mylenA\usebox\myboxB}%
    \fi}
\makeatother


%---------------------------------------------------------------------------------------------------------------
% Les dessins
\tikzstyle{every picture}+=[remember picture]
\tikzstyle{na} = [baseline=-.5ex]
\makeatletter
\newcommand{\gettikzxy}[3]{%
  \tikz@scan@one@point\pgfutil@firstofone#1\relax
  \edef#2{\the\pgf@x}%
  \edef#3{\the\pgf@y}%
}
\makeatother

% Définition des nouvelles options xmin, xmax, ymin, ymax
% Valeurs par défaut : -3, 3, -3, 3
\tikzset{
    xmin/.store in=\xmin, xmin/.default=-3, xmin=-3,
    xmax/.store in=\xmax, xmax/.default=3, xmax=3,
    ymin/.store in=\ymin, ymin/.default=-3, ymin=-3,
    ymax/.store in=\ymax, ymax/.default=3, ymax=3,
}
% Commande qui trace la grille entre (xmin,ymin) et (xmax,ymax)
\newcommand {\grille}
    {\draw[help lines] (\xmin,\ymin) grid (\xmax,\ymax);}
% Commande \axes
\newcommand {\axes}[2] {
    \draw[->,gray] (\xmin,0) -- (\xmax,0);
    \draw[->,gray] (0,\ymin) -- (0,\ymax);
    \draw (\xmax,0) node[below]{#1};
    \draw (0,\ymax) node[left]{#2};
}
\newcommand {\axesDown}[2] {
    \draw[->,gray] (\xmin,0) -- (\xmax,0);
    \draw[->,gray] (0,\ymax) -- (0,\ymin);
    \draw (\xmax,0) node[below]{#1};
    \draw (0,\ymin) node[left]{#2};
}

% Commande qui limite l’affichage à (xmin,ymin) et (xmax,ymax)
\newcommand {\fenetre}
    {\clip (\xmin,\ymin) rectangle (\xmax,\ymax);}


%---------------------------------------------------------------------------------------------------------------
% Les tableaux
\newcolumntype{C}[1]{>{\centering}p{#1}}
\newcolumntype{L}[1]{>{\raggedright}p{#1}}
\newcommand\Tstrut{\rule{0pt}{2.6ex}}         % = `top' strut
\newcommand\Bstrut{\rule[-0.9ex]{0pt}{0pt}}   % = `bottom' strut

% Tableau : epaisseur des lignes
\newcommand{\smallhrule}{\specialrule{0.02em}{0.25em}{0.25em}}
\newcommand{\medhrule}{\specialrule{0.05em}{0.3em}{0.3em}}
\newcommand{\bighrule}{\specialrule{0.1em}{0.3em}{0.5em}}


%---------------------------------------------------------------------------------------------------------------
% Math

% indicatrice
\newcommand{\ind}{\mathds{1}}

% Fonction
\newcommand{\fonction}[5]{
    \begin{array}{rcll}
        #1 \colon & #2 & \longrightarrow & #3 \\
                  & #4 & \longmapsto     & #5
    \end{array}
}

% Elements differentiels
%\newcommand{\diff}{\mathop{}\mathopen{}\mathrm{d}}
\newcommand{\xDif}{{\rm D}}
\newcommand{\xdif}{\,{\rm d}}
\newcommand{\diff}{\xdif}
\newcommand{\Diff}{\xDif}
\newcommand{\pardiff}{{\rm \partial}}
\DeclareMathOperator{\dd}{d}

% Solutions
\newcommand\rsol{\bar{r}}
\newcommand\lsol{\bar{\lambda}}
\newcommand\fsol{\bar{f}}
\newcommand\ssol{\bar{s}}
\newcommand\xsol{\bar{x}}
\newcommand\ysol{\bar{y}}
\newcommand\usol{\bar{u}}
\newcommand\zsol{\bar{z}}
\newcommand\psol{\bar{p}}
\newcommand\tfsol{\bar{t}_f}
\newcommand\tsol{\bar{t}\hspace{0.1em}}



% Notion ensembliste
\newcommand{\adherence}[1]{\xoverline{#1}}

% Les ensembles
\newcommand{\Ball}{{B}}                     % Boule ouverte
\newcommand{\BallClosed}{\xoverline{B}}     % Boule fermee

\newcommand{\M}{\mbox{$\mathbf{M}$}\xspace}
\newcommand{\U}{U}
\newcommand{\X}{X}
\newcommand{\B}{\mathcal{B}}
\newcommand{\E}{\mathcal{E}}
\newcommand{\F}{\mathcal{F}}
\newcommand{\D}{\mathcal{D}}
\renewcommand{\O}{\mathcal{O}}
\renewcommand{\P}{\mathcal{P}}

\newcommand{\Htrue}{\mathscr{H}}
\newcommand{\GLcal}{\mathscr{GL}}
\newcommand{\Ecal}{\mathscr{E}}
\newcommand{\Kcal}{\mathscr{K}}
\newcommand{\Ccal}{\mathscr{C}}
\newcommand{\Fcal}{\mathscr{F}}
\newcommand{\Ncal}{\mathscr{N}}
\newcommand{\Lcal}{\mathscr{L}}
\newcommand{\Dcal}{\mathscr{D}}
\newcommand{\Ucal}{\mathscr{U}}
\newcommand{\Vcal}{\mathcal{V}}
\newcommand{\Ical}{\mathcal{I}}
\newcommand{\Acal}{\mathscr{A}}

\newcommand{\Sgot}{\mathfrak{S}} % pour désigner le groupe des permutations ou groupe symétrique

\newcommand{\xLn}[1]{{\Lcal}^{#1}}          % Ensemble des applications lineaires continue
\newcommand{\xCn}[1]{{\Ccal}^{#1}}          % Ensemble des applications de classe C^k

% Pour le cours de mesure et integration
\newcommand{\Bor}{\mathcal{B}}
\newcommand{\AT}{\mathcal{A}}
\newcommand{\BT}{\mathcal{B}}
\newcommand{\NT}{\mathcal{N}}
\newcommand{\CT}{\mathcal{C}}
\newcommand{\OT}{\mathcal{O}}
\newcommand{\FT}{\mathcal{F}}
\newcommand{\Parties}{\mathcal{P}}
\newcommand{\tribu}{\sigma}
\DeclareMathOperator{\card}{card}
\newcommand{\FM}{\mathcal{M}} % fonction mesurable
\newcommand{\FE}{\mathcal{E}} % fonction mesurable étagée
%\newcommand{\FM}{\mathcal{F}} % fonction mesurable
%\newcommand{\FE}{\mathcal{F}^0} % fonction mesurable étagée

\newcommand{\cl}[1]{[#1]} % ou overline

\newcommand*{\EnsembleQuotient}[2]%
{\ensuremath{%
    #1\hspace{1pt}\!/\!\hspace{1pt}\raisebox{-.0ex}{\ensuremath{\mathcal{#2}}}}}

% Intervalles
\newcommand{\intervalle}[4]{\mathopen{#1}#2
                                \mathclose{}\mathpunct{},#3
                                \mathclose{#4}}
\newcommand{\intervalleff}[2]{\intervalle{[}{#1}{#2}{]}}
\newcommand{\intervalleof}[2]{\intervalle{]}{#1}{#2}{]}}
%\newcommand{\intervalleof}[2]{\intervalle{(}{#1}{#2}{]}}
\newcommand{\intervallefo}[2]{\intervalle{[}{#1}{#2}{[}}
%\newcommand{\intervallefo}[2]{\intervalle{[}{#1}{#2}{)}}
\newcommand{\intervalleoo}[2]{\intervalle{]}{#1}{#2}{[}}
%\newcommand{\intervalleoo}[2]{\intervalle{(}{#1}{#2}{)}}
\newcommand{\intervalleentier}[2]{\intervalle{\llbracket}{#1}{#2}{\rrbracket}}


% Derivees partielles
% Derivees partielles
\newcommand{\frp}[2]{\dfrac{\partial #1}{\partial #2}}
\newcommand{\frpp}[2]{\dfrac{\partial^2 #1}{\partial #2 ^2}}
\newcommand{\frpxx}[1]{\dfrac{\partial^2 #1}{\partial x^2}}
\newcommand{\frptt}[1]{\dfrac{\partial^2 #1}{\partial t^2}}
\newcommand{\frpuu}[1]{\frac{\partial^2 #1}{\partial u^2}}
\newcommand{\frpxu}[1]{\frac{\partial^2 #1}{\partial x \partial u}}
\newcommand{\frpux}[1]{\frac{\partial^2 #1}{\partial u \partial x}}
\DeclareDocumentCommand\grad{ m g }{%
    \IfNoValueTF{#2} {\nabla #1} {\nabla_{#2} #1}
}

% Ensembles de nombres
\newcommand{\nbSet}[1]{\mathbb{#1}}
\newcommand{\setPositive}{+}
\newcommand{\setNegative}{-}
\newcommand{\setStar}{\text{*}}
\newcommand{\N}{\nbSet{N}}
\newcommand{\Z}{\nbSet{Z}}
\newcommand{\Q}{\nbSet{Q}}
\newcommand{\R}{\nbSet{R}} %\newcommand{\R}{\mathbb{R}}
\newcommand{\C}{\nbSet{C}}
\newcommand{\Sn}{\nbSet{S}}
\newcommand{\K}{\nbSet{K}}
%\renewcommand{\S}{\mathbb{S}}

\newcommand{\setDeco}[2]{
    \IfEqCase{#2}{
        {b}{\xoverline{\nbSet{#1}}}
        {s}{\nbSet{#1}^{\setStar}}
        {sb}{\xoverline{\nbSet{#1}}^{\setStar}}
        {n}{\nbSet{#1}^{\phantom{\setStar}}_{\setNegative}}
        {p}{\nbSet{#1}^{\phantom{\setStar}}_{\setPositive}}
        {sn}{\nbSet{#1}^{\setStar}_{\setNegative}}
        {sp}{\nbSet{#1}^{\setStar}_{\setPositive}}
        {bp}{\xoverline{\nbSet{#1}}_{\setPositive}}
    }
}

\newcommand{\Nb}{ \ensuremath{\setDeco{N}{b}} }
\newcommand{\Ns}{ \ensuremath{\setDeco{N}{s}} }
\newcommand{\Nsb}{ \ensuremath{\setDeco{N}{sb}} }
\newcommand{\Rn}{ \ensuremath{\setDeco{R}{n}} }
\newcommand{\Rp}{ \ensuremath{\setDeco{R}{p}} }
\newcommand{\Rb}{ \ensuremath{\setDeco{R}{b}} }
\newcommand{\Rbp}{ \ensuremath{\setDeco{R}{bp}} }
\newcommand{\Rs}{ \ensuremath{\setDeco{R}{s}} }
\newcommand{\Rsp}{ \ensuremath{\setDeco{R}{sp}} }
\newcommand{\Rsn}{ \ensuremath{\setDeco{R}{sn}} }

% Matrice semi definie positive
\newcommand{\semidefpos}{\succeq 0}
\newcommand{\defpos}{\succ 0}
\newcommand{\defneg}{\prec 0}

% Valeur absolue et norme
\newcommand{\abs}[1]{\lvert#1\rvert} 
\newcommand{\absStyle}[1]{\left\lvert#1\right\rvert}
\newcommand{\norm}[1]{\lVert#1\rVert}
\newcommand{\normStyle}[1]{\left\lVert#1\right\rVert}

% Petit o et grand O
\newcommand{\petito}[1]{o\mathopen{}\left(#1\right)}
\newcommand{\grandO}[1]{O\mathopen{}\left(#1\right)}

% Ensemble des .. tels que ..
\newcommand{\enstq}[2]{\left\{#1\mathrel{}\middle|\mathrel{}#2\right\}}

% Produit scalaire
\newcommand{\prodscal}[2]{\left(#1 \,|\, #2\right)}
%\newcommand{\crochetDualite}[2]{\left\langle#1,#2\right\rangle}
\newcommand{\crochetDualite}[2]{#1 \cdot #2}
\newcommand{\Lie}[2]{\left[#1, #2\right]}

% Fleches : utiliser la commande \vv
%\renewcommand\tilde{\widetilde}

% Lettres avec tilde
\def\xt{\widetilde{x}}
\def\ut{\widetilde{u}}
\def\pt{\widetilde{p}}
\def\ft{\widetilde{f}}
\def\ct{\widetilde{c}}
\def\Et{\widetilde{E}}
\def\Ucalt{\widetilde{\Ucal}}
\def\Acalt{\widetilde{\Acal}}
\def\Phit{\widetilde{\Phi}}

% Operateur math
\makeatletter
\newcommand\RedeclareMathOperator{%
  \@ifstar{\def\rmo@s{m}\rmo@redeclare}{\def\rmo@s{o}\rmo@redeclare}%
}
% this is taken from \renew@command
\newcommand\rmo@redeclare[2]{%
  \begingroup \escapechar\m@ne\xdef\@gtempa{{\string#1}}\endgroup
  \expandafter\@ifundefined\@gtempa
     {\@latex@error{\noexpand#1undefined}\@ehc}%
     \relax
  \expandafter\rmo@declmathop\rmo@s{#1}{#2}}
% This is just \@declmathop without \@ifdefinable
\newcommand\rmo@declmathop[3]{%
  \DeclareRobustCommand{#2}{\qopname\newmcodes@#1{#3}}%
}
\@onlypreamble\RedeclareMathOperator
\makeatother

\DeclareMathOperator{\argmax}{arg\,max}
%\DeclareMathOperator{\rank}{rank}
\DeclareMathOperator{\codim}{codim}
\DeclareMathOperator{\rand}{rand}
\DeclareMathOperator{\rang}{rang}
\DeclareMathOperator{\rank}{rank}
\DeclareMathOperator{\vect}{Vect}
\DeclareMathOperator{\Isom}{Isom}
\DeclareMathOperator{\im}{Im}
\RedeclareMathOperator{\Im}{Im}
\RedeclareMathOperator{\det}{det}
\DeclareMathOperator{\trace}{tr}
\DeclareMathOperator{\comatrice}{com}
\DeclareMathOperator{\supp}{supp}
\DeclareMathOperator{\sign}{sign}
\DeclareMathOperator{\minimize}{minimiser}
\DeclareMathOperator{\comp}{num}
\DeclareMathOperator{\Ker}{Ker}
\DeclareMathOperator{\graphe}{graphe}
\RedeclareMathOperator{\ker}{Ker}
\DeclareMathOperator{\diag}{diag}
%\DeclareMathOperator{\grad}{grad}
\DeclareMathOperator{\id}{Id}
\DeclareMathOperator{\Hom}{Hom}
\DeclareMathOperator{\GL}{GL}
\DeclareMathOperator{\sh}{sh}
\DeclareMathOperator{\argsh}{argsh}
\DeclareMathOperator{\inv}{inv}
\DeclareMathOperator{\sym}{Sym}

% exponential mapping
\newcommand{\expmap}[3]{\exp({#2 #3}) (#1)}

% Fonction homotopique
\renewcommand{\hom}{\Phi}

% Epsilon, phi, theta, etc.
\newcommand{\veps}{\varepsilon}
\newcommand\vphi{\varphi}
\newcommand\pth{p_\theta}
\renewcommand\th{\theta}
\newcommand{\arc}{\gamma}

%---------------------------------------------------------------------------------------------------------------
% MACROS Ehouarn Simon
%\newcommand{\convn}{\underset{n\rightarrow +\infty}{\longrightarrow}}
%\newcommand{\convps}{\underset{n\rightarrow +\infty}{\overset{\text{p.p.}}{\longrightarrow}}}
\newcommand{\convn}{{\rightarrow}}
\newcommand{\convps}{{\overset{\text{p.p.}}{\longrightarrow}}}
